% Sources cited on the slides as (Key+YY). Full entries are in references.md.
\defref{Kosaka+24}{Kosaka, Y. et al. (2024). The JRA-3Q Reanalysis. \emph{J. Meteor. Soc. Japan}, 102, 49--109. doi:10.2151/jmsj.2024-004}
\defref{JMA+24}{Japan Meteorological Agency (2024). JRA-3Q. NSF NCAR GDEX, dataset d640000. CC BY-NC-SA 4.0}
\defref{Hersbach+20}{Hersbach, H. et al. (2020). The ERA5 global reanalysis. \emph{Q. J. R. Meteorol. Soc.}, 146, 1999--2049. doi:10.1002/qj.3803}
\defref{Gelaro+17}{Gelaro, R. et al. (2017). MERRA-2. \emph{J. Climate}, 30, 5419--5454. doi:10.1175/JCLI-D-16-0758.1}
\defref{Kobayashi+15}{Kobayashi, S. et al. (2015). The JRA-55 Reanalysis. \emph{J. Meteor. Soc. Japan}, 93, 5--48. doi:10.2151/jmsj.2015-001}
\defref{Kalnay+96}{Kalnay, E. et al. (1996). The NCEP/NCAR 40-Year Reanalysis Project. \emph{Bull. Amer. Meteor. Soc.}, 77, 437--471}
\defref{Wilks+11}{Wilks, D. S. (2011). \emph{Statistical Methods in the Atmospheric Sciences}, 3rd ed. Academic Press}
\defref{Mardia+00}{Mardia, K. V., Jupp, P. E. (2000). \emph{Directional Statistics}. Wiley}
\defref{Hastie+09}{Hastie, T., Tibshirani, R., Friedman, J. (2009). \emph{The Elements of Statistical Learning}, 2nd ed. Springer}
\defref{Harvey+96}{Harvey, V. L., Hitchman, M. H. (1996). A Climatology of the Aleutian High. \emph{J. Atmos. Sci.}, 53, 2088--2101}
\defref{Baldwin+01}{Baldwin, M. P. et al. (2001). The quasi-biennial oscillation. \emph{Rev. Geophys.}, 39, 179--229. doi:10.1029/1999RG000073}
\defref{ICAO+93}{ICAO (1993). \emph{Manual of the ICAO Standard Atmosphere}, Doc 7488/3}
